\chapter{Technische Erkundung: Proof-of-Concept mit CyberRangeCZ}
\label{chap:technische-erkundung}

Ziel der technischen Erkundung war es, die grundsätzliche Machbarkeit einer selbst betriebenen
Cyber-Range-Plattform im Sinne der späteren Masterarbeit zu überprüfen. Im Fokus stand dabei
kein umfassender Produktivbetrieb, sondern ein \emph{Proof-of-Concept} (PoC), der zentrale
Abläufe demonstriert: Installation und Deployment der Plattform, Erstzugriff über die Weboberflächen,
Import einer Trainingsdefinition, Bereitstellung einer Sandbox sowie der Zugriff auf die bereitgestellten
Systeme.

Für den PoC wurde die CyberrangeCZLite-Variante verwendet, da sie einen schnellen Einstieg in die
Plattformmechanik ermöglicht und zugleich die wesentlichen Konzepte (Portal, Trainingsdefinitionen,
Pools, Sandboxes) abbildet \cite{CyberRangeCZDevopsLite}.

\section{Testumgebung}
\label{sec:poc-testumgebung}

Der PoC wurde auf einer dedizierten Maschine mit Ubuntu 24.04 LTS durchgeführt. Die verwendeten
Ressourcen (CPU, RAM und Storage) sind in Abbildung~\ref{fig:poc-hardware} dargestellt.


\begin{figure}[H]
    \centering
    \includegraphics[width=0.85\linewidth]{poc_hardware.png}
    \caption{Testumgebung für den PoC (eigene Aufnahme)}
    \label{fig:poc-hardware}
\end{figure}

\section{Installation und Deployment der PoC-Umgebung}
\label{sec:poc-installation}

Die Installation folgt dem Ansatz, notwendige Virtualisierungs- und Container-Komponenten
lokal bereitzustellen und das eigentliche Deployment über das bereitgestellte PoC-Repository
automatisiert auszuführen. Dazu wurden zunächst Basiswerkzeuge sowie KVM/libvirt und Docker
installiert. KVM/libvirt stellt die Virtualisierungsbasis bereit (VM-Lifecycle und Netzwerk),
während Docker genutzt wird, um eine vorbereitete Umgebung auszuführen, in der eine vagrant
mit libvirt-Anbindung zur Provisionierung verwendet wird.

\begin{verbatim}
sudo apt install -y \
  git tmux curl \
  qemu-kvm libvirt-daemon-system libvirt-clients bridge-utils cpu-checker \
  docker.io
\end{verbatim}

Im nächsten Schritt wurden die relevanten Dienste aktiviert und der Benutzer für die Verwendung von
\texttt{libvirt}, KVM und Docker berechtigt. Damit wird erreicht, dass Virtualisierung und Containerbetrieb
ohne ständige Root-Kontexte nutzbar sind.

\begin{verbatim}
sudo systemctl enable --now libvirtd docker
sudo usermod -aG libvirt,kvm,docker $USER
newgrp docker
\end{verbatim}

Anschließend wurde das PoC-Repository geklont und die libvirt-URI gesetzt, sodass Vagrant eindeutig
die System-libvirt-Instanz adressiert. Zur langlebigen Durchführung lange Laufzeit, viele Ausgaben
wurde der Prozess in einer tmux-Session gestartet.

\begin{verbatim}
git clone https://github.com/cyberrangecz/devops-crczp-lite.git
cd devops-crczp-lite
export LIBVIRT_DEFAULT_URI=qemu:///system
tmux new -s crczp
\end{verbatim}

Das eigentliche Deployment erfolgt über einen Docker-Container, der Vagrant mit libvirt-Unterstützung
bereitstellt. Die relevanten Sockets und Verzeichnisse werden in den Container gemountet, damit Vagrant
(1) auf libvirt zugreifen kann, (2) seine Vagrant-States dauerhaft speichern kann und (3) im aktuellen
Arbeitsverzeichnis die Definitionsdateien findet. Der Parameter --network host stellt sicher, dass
Netzwerkoperationen und lokale Zugriffe (z.\,B. auf die Weboberfläche) ohne zusätzliche Container-NAT-
Komplexität funktionieren.

\begin{verbatim}
docker run -it --rm \
  -e LIBVIRT_DEFAULT_URI \
  -v /var/run/libvirt/:/var/run/libvirt/ \
  -v ~/.vagrant.d:/.vagrant.d \
  -v $(realpath "${PWD}"):${PWD} \
  -w "${PWD}" \
  --network host \
  vagrantlibvirt/vagrant-libvirt:latest \
  vagrant up
\end{verbatim}

Nach erfolgreichem Abschluss wurde eine URL für das Web-Portal sowie initiale Zugangsdaten ausgegeben.
Diese Ausgabe ist in Abbildung~\ref{fig:poc-deploy} dokumentiert. 


\begin{figure}[H]
    \centering
    \includegraphics[width=0.92\linewidth]{poc_deploy_done.png}
    \caption{Erfolgreiches Deployment der PoC-Umgebung (eigene Aufnahme)}
    \label{fig:poc-deploy}
\end{figure}

\section{Erstzugriff über Keycloak und CyberRangeCZ-Portal}
\label{sec:poc-webzugriff}

CyberRangeCZ verwendet für Authentifizierung und Benutzer-/Client-Verwaltung typischerweise
Keycloak als Identity- und Access-Management-Komponente. Der erfolgreiche Zugriff auf die
Keycloak-Verwaltungskonsole bestätigt, dass die Authentifizierungs- und Basisdienste funktionsfähig
bereitgestellt wurden (Abbildung~\ref{fig:poc-keycloak}).


\begin{figure}[H]
    \centering
    \includegraphics[width=0.95\linewidth]{poc_keycloak.png}
    \caption{Keycloak-Verwaltungskonsole der PoC-Installation(eigene Aufnahme)}
    \label{fig:poc-keycloak}
\end{figure}

Der zentrale Einstiegspunkt für Trainings- und Sandbox-Management ist das CyberRangeCZ Platform Portal.
Abbildung~\ref{fig:poc-portal} zeigt die Startseite und zeigt die wichtigsten Funktionen,
u.\,a.\ Trainingsdefinitionen, Trainingsinstanzen, Sandboxes, Pools sowie Benutzer- und Gruppenverwaltung.


\begin{figure}[H]
    \centering
    \includegraphics[width=0.95\linewidth]{poc_portal_home.png}
    \caption{CyberRangeCZ Platform Portal (Startseite) nach erfolgreichem Login(eigene Aufnahme)}
    \label{fig:poc-portal}
\end{figure}

\section{Import einer Trainingsdefinition und Provisionierung einer Sandbox}
\label{sec:poc-training-import}

Um die Kernfunktionalität der Plattform zu validieren, wurde eine Beispiel-Trainingsbibliothek
(\emph{library-demo-training}) aus einem Git-Repository in das Portal eingebunden \cite{CyberRangeCZLibraryDemoTraining}. Nach dem Import wurde ein Pool bzw.\ eine Trainingsinstanz erstellt und die Sandbox-Provisionierung
ausgelöst. Die resultierende Netzwerktopologie der bereitgestellten Umgebung ist in
Abbildung~\ref{fig:poc-topology} dargestellt.


\begin{figure}[H]
    \centering
    \includegraphics[width=0.80\linewidth]{poc_topology.png}
    \caption{Netzwerktopologie einer bereitgestellten Demo-Sandbox (eigene Aufnahme)}
    \label{fig:poc-topology}
\end{figure}

\section{Zugriff auf die Sandbox und funktionaler Verbindungstest}
\label{sec:poc-zugriff}

Für den Zugriff auf die bereitgestellten Systeme stellt CyberRangeCZ im PoC automatisch erzeugte
SSH-Artefakte zur Verfügung. Dazu gehören typischerweise eine Benutzer-Konfiguration sowie ein
passender privater und öffentlicher Schlüssel. Abbildung~\ref{fig:poc-ssh-files} zeigt
die entsprechenden Dateien, die nach der Sandbox-Provisionierung bereitgestellt wurden.

Neben SSH sind in Cyber-Range-Umgebungen grundsätzlich auch weitere Zugriffsmöglichkeiten relevant,
insbesondere wenn Trainings- oder Forschungsszenarien eine grafische Interaktion erfordern. Je nach
Szenariodefinition kann der Zugriff daher auch über browserbasierte Oberflächen bzw.\ GUI-Zugriffe
(z.\,B.\ über eine Web-Konsole oder Remote-Desktop-Mechanismen wie RDP/VNC/noVNC) erfolgen.
Welche Zugriffspfade tatsächlich verfügbar sind, hängt dabei von der jeweiligen Trainingsdefinition
und der bereitgestellten Infrastruktur ab.


\begin{figure}[H]
    \centering
    \includegraphics[width=0.90\linewidth]{poc_ssh_files.png}
    \caption{SSH-Daten für den Sandbox-Zugriff (eigene Aufnahme)}
    \label{fig:poc-ssh-files}
\end{figure}

Zur Verifikation, dass der Zugriff auf die Sandbox tatsächlich funktioniert und Netzwerkkommunikation
möglich ist, wurde aus der bereitgestellten Client-Umgebung ein einfacher Scan gegen das Server-System
durchgeführt. Abbildung~\ref{fig:poc-nmap} zeigt exemplarisch, dass das Zielsystem erreichbar ist und
SSH (Port 22) offen ist. Damit ist die grundlegende Nutzbarkeit der bereitgestellten Umgebung für
weitere Schritte nachgewiesen.


\begin{figure}[H]
    \centering
    \includegraphics[width=0.92\linewidth]{poc_nmap.png}
    \caption{Einfacher Funktionstest innerhalb der Sandbox (eigene Aufnahme)}
    \label{fig:poc-nmap}
\end{figure}

\section{Erkenntnisse aus dem PoC}
\label{sec:poc-erkenntnisse}

Der PoC bestätigt, dass CyberRangeCZ (in der Lite-Variante) in einer lokalen Umgebung
installiert und betrieben werden kann und dabei zentrale Plattformmechanismen bereits sichtbar werden:
versionierbare Trainingsdefinitionen, Portal-gestützte Verwaltung sowie automatisierte Provisionierung
isolierter Umgebungen. Für das Vorprojekt ist insbesondere relevant, dass sich eine Trainingsdefinition
aus einem Git-Repository importieren und anschließend reproduzierbar als Sandbox instanziieren lässt.

Für die Masterarbeit ergibt sich daraus als nächster Schritt, die technische Erkundung in Richtung
Tool-Integration zu vertiefen, d.\,h.\ konkrete Mess- und Beobachtungspunkte (Logs, Metriken,
Netzwerkdaten) in einem Use-Case gezielt zu platzieren und den Aufwand für Integration und Vergleich
mehrerer Sicherheits- bzw.\ Monitoring-Werkzeuge realistisch abzuschätzen.
